% Harvey Mudd laboratory-report starter, based on the report resources in the
% Peter N. Saeta course archive. Keep the prose and numerical examples yours.
\documentclass[11pt]{article}
\usepackage[margin=25mm]{geometry}
\usepackage{amsmath,booktabs,graphicx,xcolor}
\usepackage{fancyhdr}
\pagestyle{fancy}\fancyhf{}\lhead{Laboratory Report}\rhead{\thepage}
\setlength{\headheight}{14pt}
\newcommand{\unit}[1]{\,\mathrm{#1}}
\begin{document}
\begin{center}
{\Large\bfseries Measuring the coefficient of friction}\par\medskip
{\large Student Name \quad Physics Laboratory \quad \today}
\end{center}
\section*{Abstract}
This report summarizes an experiment, records the observations, and explains
the uncertainty in the final coefficient. A good abstract states the question,
method, numerical result, and conclusion in a few sentences.
\section{Introduction}
For a block on a horizontal surface the limiting friction force is
\(F_f=\mu_s N\). We determine \(\mu_s\) by increasing the applied force until
the block just begins to move.
\section{Experimental procedure}
Describe the apparatus clearly enough that another student could reproduce the
measurement. List calibration steps and identify the quantity read from each
instrument. Record changes to the planned procedure here.
\section{Data and analysis}
\begin{table}[h]
\centering
\caption{Example observations. Replace these values with your data.}
\begin{tabular}{ccc}
\toprule
Mass (g) & Normal force (N) & Threshold force (N)\\
\midrule
200 & 1.96 & 0.62\\
300 & 2.94 & 0.91\\
400 & 3.92 & 1.22\\
\bottomrule
\end{tabular}
\end{table}
Fit threshold force versus normal force and report the slope with its
uncertainty. Include a figure or an error-propagation calculation when useful.
\section{Discussion and conclusion}
Compare the measured value with a reference or theoretical prediction. Discuss
friction from the pulley, surface contamination, and the resolution of the
force sensor. End with one sentence stating the result and its uncertainty.
\end{document}
